% Persian print preamble. Loaded via `include-in-header` from _quarto.yml.
%
% ENGINE: LuaLaTeX, not XeLaTeX.
%
% This was settled empirically (see the Phase 0 spike). Under XeLaTeX, pandoc's
% template selects `\usepackage[bidi=default]{babel}`, which delegates to the
% `bidi` package. Its bidirectional handling is heuristic rather than a real
% implementation of the Unicode Bidirectional Algorithm, and it reverses the
% word order of Latin-script runs embedded in Persian text: "semantic line
% breaks" typeset as "breaks line semantic". For a book carrying transliterated
% Chinese and Sanskrit in most paragraphs that is disqualifying.
%
% Under LuaLaTeX the same template selects `bidi=basic`, which is a real UBA
% implementation, and the problem disappears. xepersian was also tried and
% rejected: it collides with babel's `persian` (already active via the
% `persian` documentclass option pandoc emits from `lang: fa`) and with
% pandoc's `longtable,booktabs,array` load order, and it rides the same `bidi`
% package anyway.
%
% Cost of the choice: LuaLaTeX is roughly 2-3x slower per pass than XeLaTeX,
% and Quarto runs two passes.

% Persian digits for every LaTeX counter: chapter, section, page, footnote.
% `maparabic` redefines \arabic only, so digits the translator actually typed
% are left exactly as written -- a Latin-digit year stays Latin.
\babelprovide[import, main, maparabic]{persian}

% Pandoc loads selnolig unconditionally under LuaTeX. It is an English/German
% ligature-suppression package; its Lua callback dies on Persian glyphs with
% "bad argument #1 to 'char'" and aborts the run with no `!` line in the log,
% which makes it painful to diagnose. Turn it off. The \providecommand guard
% keeps this file harmless if the engine ever changes back.
\providecommand{\selnoligoff}{}
\selnoligoff

% Footnote rule on the right. LaTeX's \footnoterule is a raw \hrule pinned to
% the left of the column and babel does not mirror it. This box inherits
% \bodydir TRT, so its *first* item lands at the right edge: rule, then \hfill.
\renewcommand{\footnoterule}{%
  \kern-3pt\hbox to \columnwidth{\vrule width .4\columnwidth height .4pt\hfill}\kern2.6pt}
